\section{Reproducibility}
\label{sec:reproducibility}

\subsection{Software and experimental configuration}

The study uses offline trajectory collection followed by supervised learning.
The experimental specification fixes benchmark membership, dimensions,
instances, optimizer seeds, portfolio composition, population size, FE budgets,
and the dynamic decision-opportunity schedule.
Reporting these elements, together with repetitions, variability, and feature
cost, is particularly important for optimization representations
\cite{cenikjSurveyFeaturesUsed2026}.

The main implementation targets Python 3.12.  The \texttt{pflacco} queries run
in an isolated Python 3.11 environment with \texttt{pflacco==1.2.2},
\texttt{numpy==1.24.3}, \texttt{pandas==2.0.3}, \texttt{scipy==1.10.1}, and
\texttt{scikit-learn==1.2.2}.  The descriptor extractor receives only the
previously evaluated LHS sample.  It neither imports a benchmark nor calls the
objective function, and a failed \texttt{pflacco} feature group is not replaced
by a study-specific substitute.

\subsection{Random streams}

All scientific random streams are constructed from explicit integers through
\texttt{numpy.random.SeedSequence}.  Optimizer streams combine the configured
optimizer seed with integer identifiers for the run and stochastic event.
Query sampling additionally uses integer stream, function, instance, dimension,
and sample-design codes.  Population-transfer initialization receives explicit
function, instance, and event codes.  The Random Analysis repetitions combine a
base seed with problem, dimension, optimizer seed, decision FE, and repetition
indices.  Experimental randomness therefore does not depend on strings or
process-specific hash values.

\subsection{Observation indexing and state alignment}

An observation is uniquely indexed by
\begin{equation}
  k_t=(\mathrm{split},p,\mathcal{F}(p),d,a_t,r,e_t),
  \label{eq:state-key}
\end{equation}
where $p$ uniquely identifies the problem instance, $\mathcal{F}(p)$ is its
function family, $a_t$ is the prefix algorithm,
$r$ is the optimizer seed, and $e_t$ is the exact integer FE count.  This index
aligns trajectories, behavior summaries, action outcomes, selection-reference
predictions, utility labels, and Decision Model observations.  Exact FE, rather
than a rounded budget ratio, identifies a shared state; the corresponding
elapsed-budget covariate is calculated as $e_t/B$.  Nominal milestone ratios and
realized native-window spans describe how the state was sampled but do not
replace exact-FE alignment.

\subsection{Data construction and evaluation dependencies}

Complete native-state trajectories and complete-budget terminal outcomes
define the common analysis population.  Permutation-invariant behavior and the
three fixed descriptor representations are evaluated at matched states, after
which the four query-adjusted actions provide the statewise losses.  Each query
configuration has its own cross-fitted selection reference, utility labels,
Decision Model, threshold, baseline comparison, and held-out evaluation.  Only
trajectories with both a complete native history and a terminal outcome enter
the corresponding estimand.  Query identity, sample design, FE charge, and
descriptor definition jointly identify an analysis condition; in particular,
the 5\% and 10\% acquisition conditions are analyzed separately and are never
pooled.

\subsection{Validity checks}

The validity checks address state continuity, representation invariance, FE
accounting, and information leakage.  For each portfolio optimizer, a restored
complete state is compared with uninterrupted native continuation at matched
checkpoints, including optimizer-specific internal dynamics rather than only
population values.  Independent reordering of population members at every
evaluated BBOB state tests that the 31 Decision Model inputs, three diagnostic
variables, and native-window summaries are permutation invariant.

Query-level checks enforce the three descriptor definitions, the shared sample
for the cheap and standard configurations, four non-duplicated actions,
continuous remaining budget, query-specific FE accounting, the statewise
action-loss transformation, and the utility and selector decompositions.
Feature-group failures remain explicit; undefined individual descriptors are
imputed only from BBOB-training observations.  Model-level checks fit all
preprocessing within the applicable training-family fold, restrict the active
candidates to LDA, Logistic Regression, and Ridge, and exclude BBOB validation
and external suites from model fitting and threshold estimation.  The three
action relations in Eq.~\eqref{eq:handoff-fields} are checked throughout
selection, utility calculation, and policy evaluation.  These procedures assess
implementation validity; they add no observations to RQ1--RQ5.

\subsection{Resource measurement}

FE accounting is exact at the state level:
\begin{align}
  FE_{\mathrm{skip,opt}} &= B-e_t,\\
  FE_{q,\mathrm{opt}} &= B-e_t-FE_q.
\end{align}
Behavior extraction reuses evaluated optimizer states and makes no additional
objective-function calls.  Query time is the sum of sample-evaluation and
feature-computation time.  Selection latency includes state preparation, model
inference, and action choice.  Prediction time is measured separately for the
T0 and B3 controllers.

All resource measurements use one fixed hardware and software environment.
Processor, operating system, Python environments, thread count, batch size,
cache state, and the peak-memory procedure accompany the estimates.  Timings
are decomposed into behavior extraction, T0 and B3 inference, query sampling,
descriptor computation, downstream selection, and total wall time.
\TBD{TBD-REPRO-RESOURCE-01}

\subsection{Recomputation of estimands}

The state-level release retains the quantities needed to recompute every term
in Eq.~\eqref{eq:query-utility} and Eq.~\eqref{eq:selector-decomposition},
including the observed losses, normalization and cost terms, selected action,
and three action relations.  Model summaries are reconstructed from
training-family fold assignments, nested and full-training OOF scores,
thresholds, and held-out predictions.  Policy summaries are reconstructed from
the trigger state, query decision, selected action, transition mode, FE use,
timing components, and terminal loss.

Equivalence is interpreted only relative to margins and procedures fixed
before policy-outcome inspection.  Held-out BBOB and each external suite are
analyzed separately with every train-derived component frozen.  These
provisions define reproducibility and do not themselves constitute evidence of
held-out or cross-benchmark transfer.
